\subsection{Web Application}

The web application is the account-management and onboarding companion to the mobile app. It is a React application deployed on Vercel and backed by a Supabase Postgres database. Sensor readings, group membership, device records, and authentication all live in Supabase; the web app and the mobile app share the same backend.

The deployed site is hosted at \url{https://cse123a-project-6a3s.vercel.app/}. The interface is organized into two side-by-side panels: an \textbf{Account} panel (the Settings page) and a \textbf{Setup guide} panel (the Guide page). The web app intentionally does not display live water-level data --- that is the role of the mobile app. Instead, the web app exists to let users manage the account that the mobile app signs in with, review the devices and groups associated with that account, and consult the first-time setup instructions.

\subsubsection{System Architecture}
The deployed architecture has three coordinated layers:

\begin{itemize}
    \item \textbf{Frontend (React + Vite):} Renders the Settings panel and the Guide panel.
    \item \textbf{Vercel backend (serverless functions):} Hosts the ingestion endpoint used by the microcontroller and the authenticated REST endpoints used by both the web app and the mobile app (e.g. \texttt{/api/groups}, \texttt{/api/devices}, \texttt{/api/app-state}).
    \item \textbf{Supabase:} Provides authentication, a Postgres database for groups/devices/readings, and row-level-security enforcement.
\end{itemize}

\subsubsection{Settings Page}
The Settings page is the Account panel of the web app. It serves two states depending on whether the user is signed in.

\paragraph{Signed out.}
The panel renders a Sign-In form with email and password fields. A hint reminds the user to sign in with the same account they use in the mobile app, since the two clients share Supabase Auth.

\paragraph{Signed in.}
After authentication, the panel displays:
\begin{itemize}
    \item The signed-in user's email and a sign-out control.
    \item \textbf{My Groups} --- a list of the groups the user belongs to, including the user's role in each group and the associated device ID where applicable. Group and device data are fetched from the authenticated endpoints \texttt{/api/groups} and \texttt{/api/devices}.
    \item \textbf{My Devices} --- a list of devices owned by the user, showing each device's name and status.
\end{itemize}

\noindent All requests from the Settings page attach the Supabase access token as a Bearer token; the serverless API routes validate the token before returning data scoped to the user.

\subsubsection{Setup Guide Page}
The Guide page is a static, step-by-step walkthrough for first-time setup of a Smart Base. It is intended to be readable on a desktop screen while the user follows along with the mobile app. The guide covers:

\begin{enumerate}
    \item Unboxing and powering on the base.
    \item Installing the mobile app.
    \item Signing in (or creating an account) and granting Bluetooth and Camera permissions.
    \item Opening the Provision Device flow from the mobile dashboard.
    \item Scanning the QR sticker on the base, which auto-fills the device name and ID.
    \item Entering the Wi-Fi SSID and password the base should join.
    \item Sending the authentication token and Wi-Fi credentials to the base over Bluetooth Low Energy.
    \item Waiting for the base to confirm Wi-Fi connectivity and begin reporting readings.
\end{enumerate}

\noindent The Guide page deliberately mirrors the mobile-app provisioning flow so users have a printable / linkable reference while completing setup on their phone.

\subsubsection{API Endpoints}
The serverless API on Vercel exposes routes used by both clients and the microcontroller:

\begin{itemize}
    \item \textbf{\texttt{POST /api/ingest}} --- Receives sensor readings (device ID, weight, battery voltage) from the microcontroller using an API key and inserts them into Supabase.
    \item \textbf{\texttt{GET /api/app-state}} --- Returns the latest reading and calibration record. Consumed by the mobile app to render the live water level.
    \item \textbf{\texttt{GET /api/groups}} --- Authenticated. Returns the groups the signed-in user belongs to. Used by both the Settings page and the mobile app.
    \item \textbf{\texttt{GET /api/devices}} --- Authenticated. Returns the devices owned by the signed-in user. Used by both the Settings page and the mobile app.
\end{itemize}

\subsubsection{Supabase Tables}
Supabase backs both the web app and the mobile app. Relevant tables include:

\begin{itemize}
    \item \textbf{\texttt{profiles}} --- One row per user. Tracks display name and an \texttt{is\_active} flag that the Settings page can toggle to deactivate an account; the mobile app subscribes to changes on this row and signs the user out when their account is deactivated.
    \item \textbf{\texttt{groups}} / \textbf{\texttt{group\_members}} --- Group definitions and role-based membership, surfaced on the Settings page under My Groups.
    \item \textbf{\texttt{devices}} --- Devices owned by a user, surfaced on the Settings page under My Devices and used during the mobile app's provisioning flow.
    \item \textbf{\texttt{water\_readings}} --- Time-series sensor uploads written by \texttt{/api/ingest}; consumed by the mobile app.
    \item \textbf{\texttt{calibration}} --- Empty/full reference points used to compute displayed percentage. Updated from the Group screen in the mobile app via \texttt{POST /api/groups/\{id\}/calibration} against the group's associated device.
\end{itemize}
